% Live-book proof adaptation of the production-verified arbitrary-field
% reduced curve packet.  Date: 28 August 2026.

\paragraph{Imported facts and the local bridge.}
Currents, positivity, multiplication of an order-zero current by a
continuous function, and differentials of currents are used as in
\cite[\S\S7.1--7.2, Example~7.2.12, and \S7.7]{CLD25}.  In particular,
\cite[Lemma~7.7.5]{CLD25} gives
\begin{equation}
\label{eq:p21p3-gfc-transpose-dd}
 \bigl\langle \mathrm d'\mathrm d''S,\varphi\bigr\rangle
 =
 \bigl\langle S,\mathrm d'\mathrm d''\varphi\bigr\rangle .
\end{equation}
Positive currents of bidimension $(0,0)$ are Radon measures by
\cite[Proposition~7.2.8]{CLD25}, and positive currents in every bidegree
have order zero by \cite[Proposition~7.2.9]{CLD25}.  We use the local
moment description of smooth forms and functions from
\cite[Proposition~15.1.1(b) and Corollary~15.1.2]{CLD25}, smooth
partitions of unity and subordinate cutoffs from
\cite[Corollary~15.2.4, Proposition~15.2.8, and
Corollary~15.2.9]{CLD25}, boundary-avoiding moments from
\cite[Lemma~12.3.6]{CLD25}, and the proper tropical pushforward in the
proof of \cite[Theorem~16.8.5]{CLD25}.  We also use
\cite[Proposition~3.3.6, Lemma~12.3.6, Theorem~12.2.4,
Corollary~12.2.5, \S14.3.2, Lemma~14.3.4,
Proposition~14.4.2, \S16.6, and Proposition~16.6.3]{CLD25}.

For an arbitrary complete ground field, the curve structure used below is
the branch--valuation correspondence, the virtual-annulus description,
and triangulation theory in
\cite[\S\S3.3.5.2, 4.2.6, 4.2.9, 4.2.10--4.2.11, and
4.3.9--4.3.10]{Ducros24CurveStructure}.  Thus an interior
type-$2$ point $x$ has a normal projective integral residual curve
$C_x$, and its branches correspond to all closed points of $C_x$.
The branch corresponding to $p$ carries the discrete valuation
$\operatorname{ord}_p$.  An interior type-$3$ point has two branches
and a virtual-annular neighborhood.  For the arbitrary-field endpoint
statement used later, \cite[Theorem~4.2.5]{Ducros24CurveStructure}
shows that an interior nonrigid point with two branches must have type
$2$ or $3$, while \cite[Theorem~4.5.4, especially \S4.5.4.1, and
\S4.5.19.1]{Ducros24CurveStructure} supplies the virtual-disc local
model and confirms that type-$1$ and type-$4$ points of a rig-smooth
curve are unibranch.  The finite-fibre and integral-affine
properties of tropical skeleta are
\cite[Corollary~3.18 and Remark~3.20]{GJR26}; finite families of
homogeneous residue functions are lifted to analytic germs by the
construction in the proof of \cite[Lemma~6.14]{GJR26}.

The proof uses \cref{lem:p21p3-gf-legal-source-localization} and exactly the
three assertions of
\cref{prop:p21p3-gf-virtual-annulus-bridge}.  In particular, it uses the
exact downstairs branch ray, finite separation of prescribed type-$2$
branches, the canonical-collar elementary star in
\cref{itm:p21p3-gf-bridge:localization}, its
global-in-\(U\) exact-preimage identity, and compatibility with the
source-separated localization of \cite[Theorem~16.8.5]{CLD25}.  It never
pulls an arbitrary current to a field extension.

\begin{proof}[Proof of \cref{lem:p21p3-gfc-rigid-zero-atom}]
The regular one-dimensional local ring $\mathcal O_{Y,y}$ has a
uniformizer represented, after shrinking, by an analytic function $f$
with $f(y)=0$.  Regard $f$ as a morphism to
$\mathbf A^{1,\mathrm{an}}_K$.  Its fibre germ at $y$ is
zero-dimensional.  Since source and target are boundaryless, the
boundary-composition identity \cite[(11.4.1.3)]{CLD25} gives
$y\in\operatorname{Int}(Y/\mathbf A^{1,\mathrm{an}}_K)$, and
\cite[Lemma~11.6.5(a)]{CLD25} makes $f$ finite at $y$.  The induced
finite homomorphism from the discrete valuation ring at $0$ to
$\mathcal O_{Y,y}$ is torsion-free and hence flat.  After shrinking as
in \cite[\S11.6.3]{CLD25}, choose a compact analytic domain $Q_{R_0}$
containing $y$, an admissible radius $R_0>0$, and a finite flat map of
constant degree $d>0$
\begin{equation}
\label{eq:p21p3-gfc-finite-rigid-tube}
             f:Q_{R_0}\longrightarrow D(0,R_0)
\end{equation}
such that $f^{-1}(0)=\{y\}$.  A finite map has empty relative boundary;
the same boundary identity puts the relative boundary of $Q_{R_0}$
over $\{|T|=R_0\}$.

Choose an admissible radius $0<R<R_0$ and fix
\[
 0<\varepsilon_0<\tfrac12\log(R_0/R).
\]
For an admissible $0<r<R$ with
$\log(R/r)>4\varepsilon_0$, put $s_R=-\log R$ and
$s_r=-\log r$.  Choose a $C^\infty$ function
$\rho_r:\mathbf R\to[0,1]$ which is zero on
\[
 (-\infty,s_R-\varepsilon_0]\cup[s_r+\varepsilon_0,\infty),
\]
is one on $[s_R,s_r]$, is nondecreasing on
$[s_R-\varepsilon_0,s_R]$, and is nonincreasing on
$[s_r,s_r+\varepsilon_0]$.  Set
\[
 q_r(t):=\int_{-\infty}^{t}\rho_r(\tau)\,\mathrm d\tau,
 \qquad
 H_r:=\int_{\mathbf R}\rho_r(\tau)\,\mathrm d\tau .
\]
Then
\begin{equation}
\label{eq:p21p3-gfc-radial-data}
 H_r\geq\log(R/r),
 \qquad
 \int_{\mathbf R}|q_r''(t)|\,\mathrm dt=2.
\end{equation}
On the punctured tube put
$\chi_r=q_r(-\log|f|)$.  Extend it by zero across the outer region and
by the constant $H_r$ near $y$.  The extensions occur where $q_r$
is constant, so $\chi_r$ is a compactly supported CLD-smooth function,
with support in the fixed compact domain $Q_{R_0}$.

On each of the two transition tubes, the map $s=-\log|f|$ projects onto
the corresponding target interval and
\begin{equation}
\label{eq:p21p3-gfc-radial-dd}
 \mathrm d'\mathrm d''\chi_r
   =q_r''(s)\,\mathrm d's\wedge\mathrm d''s .
\end{equation}
The sign of $q_r''$ is fixed on each transition band.  Restrict proper
calibrated pushforward under
\eqref{eq:p21p3-gfc-finite-rigid-tube} to either transition tube.  By
\cite[Proposition~14.4.2]{CLD25} it is the corresponding target annulus
with total local degree $d$; in particular all sheets and residue-degree
factors are included.  If $\mu_f$ and $\mu_D$ denote the source and radial
calibrated measures, then $f_*\mu_f=d\mu_D$ on each transition band.
Since the two bands have disjoint supports and $q_r''$ has one sign on
each of them, \eqref{eq:p21p3-gfc-radial-data} gives
\begin{equation}
\label{eq:p21p3-gfc-radial-tv}
 \left\|\mathrm d'\mathrm d''\chi_r\wedge\delta_Y\right\|_{\mathrm{TV}}
 =d\int_{\mathbf R}|q_r''(t)|\,\mathrm dt=2d.
\end{equation}

The measure $\mu_u=\mathrm d'\mathrm d''[u]$ is positive by
\cite[Lemma~8.5.5]{CLD25}.  Positivity,
\eqref{eq:p21p3-gfc-transpose-dd}, and
\eqref{eq:p21p3-gfc-radial-tv} yield
\begin{align}
 \log(R/r)\,\mu_u(\{y\})
 &\leq \int_Y\chi_r\,\mathrm d\mu_u                                      \notag\\
 &=\bigl\langle[u],\mathrm d'\mathrm d''\chi_r\bigr\rangle                \notag\\
 &\leq 2d\sup_{Q_{R_0}}|u|.
\label{eq:p21p3-gfc-radial-bound}
\end{align}
The right side is independent of $r$.  Nontriviality of the valuation
supplies admissible radii tending to zero, so
\eqref{eq:p21p3-gfc-radial-bound} proves $\mu_u(\{y\})=0$.
\end{proof}

\begin{proof}[Proof of \cref{prop:p21p3-gfc-smooth-rigidity}]
We first record three ingredients and then globalize them.

\paragraph{The Euclidean calculation.}
Let $\Gamma$ be a locally finite one-dimensional polyhedral set in an
open subset of $\mathbf R^N$, and let $S$ be a positive current of
bidimension $(1,1)$, supported on $\Gamma$, with
$\mathrm d'S=\mathrm d''S=0$.  Positivity and
\cite[Proposition~7.2.9]{CLD25} write $S$ as a symmetric positive
semidefinite matrix $M=(M_{jk})$ of Radon measures.  Closedness is the
distributional system
\begin{equation}
\label{eq:p21p3-gfc-matrix-divergence}
                   \sum_j\partial_jM_{jk}=0
                   \quad(1\leq k\leq N).
\end{equation}
On the relative interior of an edge, choose affine coordinates
$(t,z_2,\ldots,z_N)$ in which the edge is
$\{z_2=\cdots=z_N=0\}$.  Testing
\eqref{eq:p21p3-gfc-matrix-divergence} against
$\psi(t)z_\ell$, for $\ell>1$, kills every transverse row and column;
testing against $\psi(t)$ then makes the remaining measure a constant
multiple of Lebesgue measure.  Thus
\begin{equation}
\label{eq:p21p3-gfc-edge-density}
                     S|_e=b_e[\,e\,],
                     \qquad b_e\geq0,
\end{equation}
with $b_e$ constant.  Endpoint terms at a vertex give
\begin{equation}
\label{eq:p21p3-gfc-balancing}
                       \sum_{e\ni v}b_eu_{e,v}=0,
\end{equation}
where $u_{e,v}$ is the primitive outgoing tangent.  A positive remainder
at $v$ would have the form $P\delta_v$ with $P$ positive
semidefinite.  Once the endpoint terms balance, tests vanishing at $v$
with arbitrary first jet isolate the derivatives of $P\delta_v$ in
\eqref{eq:p21p3-gfc-matrix-divergence}; hence $P=0$.  In the vertex
applications below,
\cref{itm:p21p3-gf-bridge:localization} produces a one-vertex source
neighborhood on which the full nonconstant target star of the fixed finite
moment is finite.  We apply \eqref{eq:p21p3-gfc-balancing} to that entire
finite target star.  We never infer a balancing identity by deleting rays
and passing to an arbitrary finite substar.

\paragraph{Legitimate source-separated localization.}
We make explicit why no smooth plateau is being treated as an analytic
moment coordinate.  Let $\alpha$ be a compactly supported smooth
$(1,1)$-test form.  Start by localizing in the source.  The local moment
description \cite[Proposition~15.1.1(b)]{CLD25}, a smooth partition of
unity, compactness of $\operatorname{Supp}(\alpha)$, the finite-fibre
assertion \cite[Corollary~3.18]{GJR26}, and subdivision give finitely many
summands, each contained either in one relatively compact
virtual-annular collar or in one star neighborhood of a retained vertex.
The two collars at a type-$3$ point may be retained in its single star.
An edge neighborhood meets no other retained source edge, and a star
neighborhood contains no other retained vertex.  This source-first
operation is made before any Euclidean pushforward; distinct edge pieces
are never put into one compact source domain merely in order to separate
them by a cutoff.

Fix one summand in a source neighborhood $W$.  On an edge collar, apply
\cref{lem:p21p3-gf-legal-source-localization} directly.  At a type-$3$
center, first include in its prescribed block a $K$-analytic unit with
nonzero annular slope and shorten the two collars; the existence and
two-sided monomial norm formula are part of
\cite[\S\S4.2.11.2 and 4.3.9--4.3.10]{Ducros24CurveStructure}.  Apply the
same lemma after this shortening.

Suppose now that $W$ is a one-vertex star about a type-$2$ point $x$.
For the comparison at hand, let $A$ be the finite family of branches that
must be isolated; it contains the branches on which the form block is
nonconstant and, in the divisor calculation below, the finite support of
the principal divisor being tested.  We invoke
\cref{itm:p21p3-gf-bridge:localization} with $T$, this exact test summand
$\alpha$, this exact $W$, and the selected family $A$, and spell out its
two-stage output.  Its first application of
\cref{lem:p21p3-gf-legal-source-localization} produces a legal form and
retained-cutoff moment $F_0$.  After the branch block is appended, write
$G_0=(F_0,F_{\mathrm{br}})$.  The bridge then shrinks to its canonical-collar
one-vertex source neighborhood $W_1\subset W$.  Every virtual-disc
component is constant for $G_{0,\mathrm{trop}}$ on its entirety; only
finitely many retained virtual-annular components can be active, and every
selected active ray is separated from every other active branch.

Crucially, the bridge next reapplies
\cref{lem:p21p3-gf-legal-source-localization} \emph{inside $W_1$}, with
$G_0$ as its prescribed block.  We use its re-localized output
\begin{equation}
\label{eq:p21p3-gfc-relocalized-output}
 V\subset W_1,\qquad F,\qquad
 \Omega=\mathbf R^N\setminus
 F_{\mathrm{trop}}\bigl(\partial(V/W_1)\bigr),\qquad
 U=F_{\mathrm{trop}}^{-1}(\Omega).
\end{equation}
For every selected branch $b_i\in A$, it supplies the full canonical collar
$Z_{b_i}\subset U$ and the exact identity
\begin{equation}
\label{eq:p21p3-gfc-imported-exact-preimage}
 F_{\mathrm{trop}}^{-1}\!\left(
 F_{\mathrm{trop}}(x)+(0,\varepsilon_i)v_i\right)\cap U
 =Z_{b_i},
\end{equation}
which is precisely
\eqref{eq:p21p3-gf-source-separated-ray-preimage}.  Thus the isolated
target ray receives exactly its named source collar, not merely a branch
with the same ray germ.  This is the source-separation fact used below.

The full moment $F$ has finitely many analytic-unit coordinates.  By
\cite[Proposition~4.2.4(ii)]{Ducros24CurveStructure}, only finitely many
branches at $x$ are nonconstant for at least one of them.  Hence the
\emph{full} nonconstant target star in this localized moment is finite.
All of its rays, including rays activated by the second localization, are
retained when we apply Euclidean balancing.  We do not discard the
remaining geometric branches by a finite-fibre argument: they are constant
for $F$ on cofinal collars and map to the vertex.

In every case, the legal source-localization lemma gives a genuine analytic
block $H$ and a smooth $\Phi$ with
$\chi=\Phi\circ H_{\mathrm{trop}}$ whenever a source cutoff $\chi$ is
retained.  We append $H$, never $\chi$.  It also supplies the compact
domain $V$, the full moment $F$, and the proper map
$F_{\mathrm{trop}}|_U:U\to\Omega$; in the type-$2$ case these are exactly
the re-localized objects in
\eqref{eq:p21p3-gfc-relocalized-output}.  Thus no nonlinear cutoff is
declared to be a moment coordinate, and the current is never multiplied by
a cutoff.  The proper Euclidean pushforward
\begin{equation}
\label{eq:p21p3-gfc-euclidean-pushforward}
 \langle T_\Omega,\beta\rangle
       :=\langle T|_U,F^*\beta\rangle,
 \qquad \beta\in A_c^{1,1}(\Omega),
\end{equation}
defines a positive closed Euclidean current supported on a locally finite
one-dimensional polyhedral set.  Restriction to the open set $U$, not
multiplication by a cutoff, preserves closedness.  In a type-$2$ star,
\eqref{eq:p21p3-gfc-imported-exact-preimage} ensures that an isolated
selected ray receives the contribution of exactly its named source collar.

On such a collar, let $b_e\geq0$ be the flux of $T_\Omega$ from
\eqref{eq:p21p3-gfc-edge-density}, and let $m_e>0$ be the corresponding
flux for the canonical integration current.  Put
\begin{equation}
\label{eq:p21p3-gfc-branch-ratio}
                              a_e:=b_e/m_e.
\end{equation}
On a common source-separated concatenation of two moments, tropical
pushforward multiplies $b_e$ by the tangent-lattice expansion index:
the tangent tensor supplies its square and one-dimensional measure change
supplies its inverse.  The calibrated flux $m_e$ is multiplied by the
same index by \cite[\S14.3.2, Lemma~14.3.4, and
Proposition~14.4.2]{CLD25}.  Therefore $a_e$ is independent of the
moment, subdivision, and localized presentation, and is constant along
the virtual-annular collar.  This comparison is entirely downstairs over
$K$.

\paragraph{Comparison at type-$3$ and type-$2$ points.}
At an interior type-$3$ point the two branch rays in the refined moment
are opposite rays of one line.  Applying
\eqref{eq:p21p3-gfc-balancing} first to $T$ and then to $\delta_Y$
shows that the ratios \eqref{eq:p21p3-gfc-branch-ratio} on the two sides
are equal.  Possible virtual-annulus expansion factors occur equally in
the ordinary and calibrated fluxes and have already cancelled in the
ratio.

Let $x$ now be an interior type-$2$ point, and let $C_x$ be its
normal projective integral residual curve.  For a closed point
$p\in C_x$, use a rational function $s_p$ with
$\operatorname{ord}_p(s_p)=1$, a norm-one analytic lift, and
\cref{itm:p21p3-gf-bridge:ray}.  Apply
\cref{itm:p21p3-gf-bridge:localization} with the $p$-branch among the
selected branches.  Its re-localized exact-preimage identity
\eqref{eq:p21p3-gfc-imported-exact-preimage} makes the resulting ray the
image of precisely that source collar.  The outgoing downstairs vector
has a coordinate equal to one and is therefore primitive.  Denote the
calibrated flux on this isolated branch ray by $m_p>0$, and its ratio by
$a_p$.  These quantities are independent of the normalized moment: in
a common concatenation the source tangent and both projected tangents
retain a coordinate equal to one, so both projection indices are one.

Define additive maps on the divisor group of $C_x$ by
\begin{equation}
\label{eq:p21p3-gfc-divisor-functionals}
 L_T([p])=a_pm_p,
 \qquad
 L_Y([p])=m_p.
\end{equation}
They are nonnegative on effective divisors.  If
$g\in\operatorname{Rat}(C_x)^\times$, append a norm-one analytic lift
of $g$, and apply the source-separated localization in
\cref{itm:p21p3-gf-bridge:localization} simultaneously to the finite set
$\operatorname{Supp}(\operatorname{div}(g))$.  In the scalar coordinate
of this lift, the primitive slope on the $p$-branch is
$\operatorname{ord}_p(g)$; every branch outside that finite support has
slope zero.  The scalar component of
\eqref{eq:p21p3-gfc-balancing} gives
\begin{equation}
\label{eq:p21p3-gfc-T-principal}
       0=\sum_p a_pm_p\operatorname{ord}_p(g)
        =L_T(\operatorname{div}(g)).
\end{equation}
The integration current is closed by \cite[\S7.7.4]{CLD25}; applying the
same calculation to $\delta_Y$ gives
\begin{equation}
\label{eq:p21p3-gfc-Y-principal}
       0=\sum_p m_p\operatorname{ord}_p(g)
        =L_Y(\operatorname{div}(g)).
\end{equation}
No infinite sum and no cutoff on an infinite-valence star is used.

We use the following elementary consequence of Riemann--Roch over the
actual residue field.  Let $C$ be a normal projective integral curve over
a field $k$, and let
$L:\operatorname{Div}(C)\to\mathbf R$ be additive, nonnegative on
effective divisors, and zero on principal divisors.  If $D$ has degree
zero, choose an effective divisor $A$ such that
$\deg A+\chi(C,\mathcal O_C)>0$.
For every $n\in\mathbf Z$, Riemann--Roch gives
\[
 h^0(C,\mathcal O_C(A+nD))
 \geq \deg A+\chi(C,\mathcal O_C)>0.
\]
Thus $A+nD$ is linearly equivalent to an effective divisor and
\[
                         0\leq L(A)+nL(D)
                         \qquad(n\in\mathbf Z).
\]
Letting $n$ tend to both signs of infinity gives $L(D)=0$.  Choose
any closed point $p_0$ of $C$.  Since
$\deg(p_0)D-\deg(D)[p_0]$
has degree zero,
\begin{equation}
\label{eq:p21p3-gfc-degree-functional}
                 L(D)=c_L\deg(D),
 \qquad
                 c_L:=\frac{L([p_0])}{\deg(p_0)}\geq0.
\end{equation}
This argument neither assumes geometric integrality nor uses a degree-one
closed point.

Apply \eqref{eq:p21p3-gfc-degree-functional} to $L_Y$ and $L_T$.
There are $c_Y>0$ and $c_T\geq0$ such that
\[
                  m_p=c_Y\deg(p),
 \qquad
                  a_pm_p=c_T\deg(p)
\]
for every closed point $p$.  Hence
\begin{equation}
\label{eq:p21p3-gfc-equal-star-ratios}
                              a_p=c_T/c_Y
\end{equation}
on every branch at $x$.

\paragraph{Globalization on the smooth curve.}
Let $E\subset Y$ be the set of type-$1$ and type-$4$ points.  We justify
the endpoint assertion over the original field.  At a nonrigid point,
\cite[Theorem~4.2.5]{Ducros24CurveStructure} says that the existence of
two branches forces type $2$ or $3$; purity and the local graph structure
give at least one branch.  At a rigid point the virtual-disc neighborhood
of \cite[Theorem~4.5.4 and \S4.5.4.1]{Ducros24CurveStructure} gives the
same conclusion.  Thus every point of $E$ is unibranch, as is also
recorded for nonrigid type-$1$/type-$4$ points in
\cite[\S4.5.19.1]{Ducros24CurveStructure}.

Connected analytic spaces are arcwise connected by
\cite[\S12.1.12]{CLD25}.  If two points lie outside $E$, choose an arc
joining them.  By \cite[\S1.3.12]{Ducros24CurveStructure}, a point lying
in the open interval of a graph has at least the two branches supplied by
the two sides.  Hence no unibranch point of $E$ can occur in the interior
of the chosen arc, and $Y\setminus E$ is connected.  The ratio is
constant on each annular collar, propagates through type-$3$ points by
the two-ray comparison, and propagates through type-$2$ points by
\eqref{eq:p21p3-gfc-equal-star-ratios}.  It is therefore one constant
$c\geq0$ on $Y\setminus E$.

The source-separated local moments tropicalize every local test form.
The edge-density calculation and the absence of a vertex remainder then
show $T=c\delta_Y$ on a neighborhood of each point of $Y\setminus E$.
Every point of $E$ has tropical dimension zero, so
\cite[Proposition~3.3.6]{CLD25} says that the germ of every smooth
$(1,1)$-test form vanishes there.  Thus the support of each such test
form is disjoint from $E$, and the current $T-c\delta_Y$, already zero
off $E$, pairs to zero with every test form.  This proves
\eqref{eq:p21p3-gfc-smooth-rigidity}.  Uniqueness follows by testing on a
nonempty virtual-annular band.  All arguments were made on compact test
supports, so $Y$ need not be compact.
\end{proof}

\begin{proof}[Proof of \cref{prop:p21p3-curve-rigidity}]
We first pass from the smooth normalizations to the reduced curve.  On a
good affinoid domain write
\[
 \mathcal Q:=\nu_*\mathcal O_{\widetilde X}/\mathcal O_X,
 \qquad
 \mathfrak c:=\operatorname{Ann}_{\mathcal O_X}(\mathcal Q),
\]
and let $Z$ be the union of
\[
 \operatorname{Supp}(\mathcal Q)
   =V(\mathfrak c)
   =\operatorname{Supp}(\mathcal O_X/\mathfrak c)
\]
with the pairwise intersections of distinct irreducible components.
Analytic local rings on a good space are excellent by
\cite[\S11.1.10]{CLD25}; normalization is finite and its components
correspond to the irreducible components by
\cite[\S11.3.1]{CLD25}.  Hence $\mathcal Q$ and the conductor are
coherent.  Their support misses every generic point of the reduced pure
curve and is therefore zero-dimensional.  The same dimension argument
applies to proper intersections of distinct components.  Irreducible
components are locally finite by \cite[\S11.2.5]{CLD25}; consequently
$Z$ is closed and locally finite.

Strictness is used here.  A zero-dimensional strict affinoid algebra is
Artinian, and the strict affinoid Nullstellensatz makes the completed
residue field of each of its points finite over $K$.  Thus every point
of $Z$ is rigid.  A finite morphism induces a finite extension of
completed residue fields, so every point of
\begin{equation}
\label{eq:p21p3-gfc-conductor-preimage}
                         Z_i:=\nu_i^{-1}(Z),
\end{equation}
is rigid as well, though it need not be $K$-rational.  Moreover, a finite
algebra over a strict affinoid algebra is strict.  Since normalization is
finite, every $\widetilde X_i$ is therefore strict.

The finite normalization map is topologically proper, has empty relative
boundary, and is an isomorphism away from $Z_i$.  Applying the
boundary-composition identity \cite[(11.4.1.3)]{CLD25} to
$X_i\to X\to\mathcal M(K)$ and then to
$\widetilde X_i\to X_i\to\mathcal M(K)$ shows that
$\widetilde X_i$ is boundaryless.  It is connected because it is the
normalization of one irreducible component.

The fundamental-cycle identity is
\begin{equation}
\label{eq:p21p3-gfc-fundamental-cycle}
             \delta_X=\sum_i(\nu_i)_*\delta_{\widetilde X_i}.
\end{equation}
Indeed, a reduced pure curve has generic multiplicity one on every
irreducible component.  Integration agrees with the associated cycle by
\cite[Proposition~16.6.3]{CLD25}, and normalization has calibrated degree
one over the normal locus.  The complement is the locally finite rigid
set $Z$, where every $(1,1)$-test-form germ vanishes by
\cite[Proposition~3.3.6]{CLD25}; this proves
\eqref{eq:p21p3-gfc-fundamental-cycle} everywhere.

The complement $\widetilde X_i\setminus Z_i$ is open, good, separated,
strict, rig-smooth, and boundaryless.  Deleting $Z_i$ does not disconnect
$\widetilde X_i$.  Indeed, two
remaining points can be joined by an arc, and
\cite[\S1.3.12]{Ducros24CurveStructure} shows that no rigid unibranch
point can lie in the interior of that arc.  On $X\setminus Z$, the component
images are disjoint and normalization is an analytic isomorphism.
Applying \cref{prop:p21p3-gfc-smooth-rigidity} to each connected
$\widetilde X_i\setminus Z_i$ gives
\[
                         T=\sum_i c_i\delta_{X_i},
                         \qquad c_i\geq0,
\]
there.  Set $S=T-\sum_i c_i\delta_{X_i}$.  It has zero germ off $Z$,
while every $(1,1)$-test form has support disjoint from $Z$ by
\cite[Proposition~3.3.6]{CLD25}.  Hence $S=0$, proving
\eqref{eq:p21p3-curve-classification}.  The sum is locally finite,
and uniqueness follows by restricting to an open normal locus of $X_i$
meeting no other component.

It remains to prove the Bedford--Taylor assertion.  Put
$\widetilde W_i=\nu_i^{-1}(W)$.  Multiplying the order-zero identity
\eqref{eq:p21p3-gfc-fundamental-cycle} by a continuous function $F$ on
$W$, using \cite[Example~7.2.12]{CLD25}, gives
\begin{equation}
\label{eq:p21p3-gfc-function-cycle-decomposition}
 [F]=F\delta_W
     =\sum_i(\nu_i)_*
        \bigl((\nu_i^*F)\delta_{\widetilde W_i}\bigr).
\end{equation}
The projection identity used here follows directly by testing against a
smooth form, and the sum is finite on each test support.

Pullback preserves continuous plurisubharmonic functions by
\cite[Proposition~8.1.5]{CLD25}.  On $\widetilde W_i$, set
\[
 \mu_i^U=\mathrm d'\mathrm d''[\nu_i^*U],
 \qquad
 \mu_i^V=\mathrm d'\mathrm d''[\nu_i^*V].
\]
These are positive Radon measures by \cite[Lemma~8.5.5]{CLD25}.
Proper pushforward commutes with the differential on currents by
transposition; equivalently, use \cite[Proposition~8.5.6]{CLD25}.
Applying \eqref{eq:p21p3-gfc-function-cycle-decomposition} to $U$ and
$V$, the curvature equality
\eqref{eq:p21p3-curve-equal-curvature} becomes
\begin{equation}
\label{eq:p21p3-gfc-component-measure-sum}
                    \sum_i(\nu_i)_*(\mu_i^U-\mu_i^V)=0.
\end{equation}

Away from $Z$, the component images are disjoint and $\nu_i$ is an
isomorphism.  Hence each signed measure $\mu_i^U-\mu_i^V$ is supported
on the locally finite rigid set $Z_i\cap\widetilde W_i$.  At any point of
this set, apply \cref{lem:p21p3-gfc-rigid-zero-atom} on a connected open
neighborhood contained in $\widetilde W_i$.  Neither positive measure has
an atom there.  A signed Radon measure supported on a locally finite set
is the locally finite sum of its point masses.  Therefore
\begin{equation}
\label{eq:p21p3-gfc-component-measure-equality}
                              \mu_i^U=\mu_i^V
                              \quad\text{for every }i.
\end{equation}

Finally put $h=U-V$.  By
\eqref{eq:p21p3-curve-classification},
\[
 hT|_W=\sum_i c_i(\nu_i)_*
       \bigl((\nu_i^*h)\delta_{\widetilde W_i}\bigr).
\]
Proper pushforward commutes with $\mathrm d'\mathrm d''$, and
\eqref{eq:p21p3-gfc-component-measure-equality} gives
\begin{align*}
 \mathrm d'\mathrm d''(hT|_W)
 &=\sum_i c_i(\nu_i)_*
       \mathrm d'\mathrm d''[\nu_i^*U-\nu_i^*V] \\
 &=\sum_i c_i(\nu_i)_*(\mu_i^U-\mu_i^V)=0.
\end{align*}
This is \eqref{eq:p21p3-curve-bt-conclusion}.  It proves the
independence condition of \cite[Definition~8.6.2]{CLD25} for the
continuous psh pairs stated here; no claim about discontinuous psh
functions is made.
\end{proof}
